\answer{ex:06:1}{$e=(1-e^{-T_a/\tau_e})e^{-d/\tau_e}$. (a)~$0.110$ frente a $4.5\cdot10^{-5}$, unas $\approx2.5\cdot10^3$ veces. (b)~$\tau_e^*=T_a/\ln(1+T_a/d)=0.59\,\text{s}$.}
\answer{ex:06:2}{Tasa $\nu_x\nu_y(A_+\tau_+-A_-\tau_-)=0.8A_+$ por segundo, $\approx2900A_+$ por hora; $\tau_-=\tau_+/0.6=33\ms$.}
\answer{ex:06:3}{$\mathbf e_1$ si $\mu^2+s_1^2>s_2^2$; aquí $1.01>0.25$, y la neurona aprende $\mathbf e_1$, aunque la dispersión a lo largo de $\mathbf e_2$ es $25$ veces mayor: la regla halla el vector principal de la matriz de correlaciones, no de covarianzas.}
\answer{ex:06:4}{$V=Re^{-(t_c+L-t)/\tau_\gamma}$ entre la lámpara y el jugo, así que $\dot V=V/\tau_\gamma$; pico de área $Re^{-L/\tau_\gamma}$: $0.61R$ y $0.78R$.}
\answer{ex:06:5}{Relación señal-ruido $1/(N+1)$, intentos $\propto N+1$: $5\cdot10^5$ intentos $\approx1$~mes y $5\cdot10^{11}$ intentos $\approx8\cdot10^4$~años.}
\answer{ex:06:6}{(a)~$k_2=1/\tau_d$, $k_1=1/\tau_d-1/\tau_\gamma$. (b)~Error falso $-(V/\tau_\gamma)\,\varepsilon/(1+\varepsilon)$, $\varepsilon=\tau_d/\tau_\gamma$, de donde $\tau_d\le0.105\,\text{s}$.}
\answer{ex:06:7}{$0.46$, $0.90$, $0.99$, $0.95$ y $0.70$ con $d=3\,\text{s}$. Los puntos caen en una sola curva de la fracción de traza $F=(1-e^{-T_{\text{cue}}/\tau_e})e^{-d/\tau_e}$: aprendizaje con $F\gtrsim0.1$, elección al azar con $F<10^{-2}$.}
\answer{ex:06:8}{$\eta^*=1/\bigl(2\sigma^2(N+2)\bigr)$, en $N+2$ intentos; con $m$ fluctuando de $N$, en $N(m+2)/m\ge N+2$: la búsqueda dispersa no ayuda.}
